% !TEX root = ./main.tex
\chapter{Conclusioni e Sviluppi Futuri}
\label{chap:conclusioni}

\section{Bilancio del lavoro e contributi principali}
\label{sec:sintesi}
\label{sec:contributi-raggiunti}
Il lavoro parte da un problema concreto: far funzionare una conversazione vocale credibile dentro un ambiente 3D WebGL, senza ridurre l'avatar a un semplice ``parlante'' scollegato dalla scena. Da qui deriva anche l'impostazione del prototipo: l'avatar viene trattato come un profilo operativo composto da aspetto, voce e memoria, sostenuto da una pipeline STT--RAG/LLM--TTS e da una separazione chiara tra frontend Unity e backend a micro-servizi. Durante lo sviluppo molte scelte non sono nate in astratto, ma dall'integrazione concreta del sistema. Anche la selezione finale dei modelli (Whisper, \texttt{llama3:8b-instruct-q4\_K\_M} in locale e \texttt{llama3.1:8b} nell'ambiente Ubuntu usato per le prove, Coqui XTTS v2, \texttt{nomic-embed-text}) è arrivata dopo tentativi reali, con cambi di rotta spesso costosi in termini di tempo. Le misure quantitative disponibili e i relativi limiti metodologici sono raccolti nel Capitolo~\ref{chap:risultati}; qui interessa soprattutto fissare cosa ha retto davvero e cosa invece resta ancora da consolidare.

Il risultato più solido è la tenuta della pipeline end-to-end, abbastanza leggibile da permettere diagnosi rapide e abbastanza stabile da funzionare sia in locale sia dietro reverse proxy. La normalizzazione degli endpoint lato client è stata decisiva, perché ha permesso di passare da loopback a dominio pubblico senza cambiare logica applicativa. Anche sul piano della riproducibilità il prototipo resta ispezionabile in modo diretto: il codice sorgente sviluppato per SOULFRAME è disponibile nella repository del progetto \cite{SOULFRAMERepo}.

Sul frontend, la FSM unica per desktop e touch, affiancata da \texttt{UINavigator} e \texttt{UIHintBar}, ha mantenuto coerente il comportamento dell'interfaccia al variare degli input. Nella conversazione, streaming TTS e \emph{wait phrases} non eliminano i tempi di inferenza, ma riducono i vuoti percepiti e rendono più leggibile lo stato del turno. Un altro punto riuscito riguarda la gestione avatar nei vincoli di WebGL: delegare import, catalogo e serving al backend ha ridotto la fragilità lato browser proprio nei punti in cui cache, persistenza e filesystem virtuale tendevano a divergere. Nel complesso il prototipo riesce a tenere insieme credibilità del turno conversazionale e continuità del profilo avatar, mostrando anche un compromesso chiaro: la build Windows resta più favorevole sul piano della continuità audio, mentre WebGL è la scelta più sensata quando contano accesso immediato e distribuzione semplificata.

\section{Limiti attuali del sistema}
\label{sec:limiti}
Il limite metodologico principale è già comparso nel Capitolo~\ref{chap:risultati}: i tempi riportati sono ordini di grandezza osservati, non misure turn-based strumentate, perché manca una telemetria persistente e sincronizzata tra client e servizi. L'impatto pratico è stato chiaro fin dalle prime prove: senza timestamp per turno la diagnosi aiuta il debugging operativo, ma non regge da sola una valutazione comparativa robusta.

Un nodo ancora aperto riguarda la memoria conversazionale. Oggi SOULFRAME memorizza in modo affidabile quando c'è un trigger lessicale esplicito (per esempio ``ricorda che...''), ma non estrae ancora in autonomia le informazioni davvero memorabili dal dialogo libero. Inoltre il contesto recente di sessione resta confinato alla RAM e al log del turno: aiuta a mantenere continuità locale, ma non viene promosso automaticamente nella memoria persistente dell'avatar. Per ora si è preferito non forzare una versione ``semi-intelligente'', perché il rischio di falsi positivi in memoria sarebbe stato più dannoso del beneficio atteso. Questo punto si lega al perimetro complessivo del sistema: SOULFRAME rimane pensato per un'interazione individuale per sessione, non per una conversazione condivisa e sincronizzata tra più utenti o più agenti concorrenti.

Un altro limite riguarda la personalizzazione facciale spinta a partire da immagini reali. I primi tentativi con pipeline di ricostruzione più autonome hanno mostrato presto che il costo di integrazione era troppo alto rispetto al beneficio visivo ottenibile nel perimetro attuale; per questo questa direzione non è ancora entrata in un flusso operativo abbastanza stabile da affiancare il percorso avatar già adottato. Avaturn resta quindi la soluzione più solida per mantenere qualità visiva coerente e complessità implementativa sostenibile.

Restano poi i vincoli del runtime WebGL: i test operativi si sono concentrati su Microsoft Edge (Chromium), mentre altri browser non sono stati caratterizzati in modo sistematico. Nel progetto sono presenti anche \texttt{fs.jslib} e l'uso di \texttt{FS.syncfs} come supporto al filesystem virtuale del browser, ma nel flusso corrente del catalogo avatar WebGL la fonte principale è il backend; per questo eventuali stati intermedi del filesystem client incidono meno sulla disponibilità degli avatar importati rispetto a una persistenza puramente browser-side. Il compromesso è chiaro: meno autonomia locale del browser, ma maggiore consistenza del catalogo quando la stessa build deve funzionare sia in test locale sia dietro proxy pubblico. Nella stessa area rientra anche il comportamento audio: il costo della mediazione WebGL incide più sulla continuità del parlato che sulla correttezza funzionale del turno. La build Windows, provata sullo stesso impianto applicativo, ha mostrato una resa più stabile e meno soggetta a micro-interruzioni percepite. Qui il compromesso progettuale si vede bene: scegliere WebGL semplifica accesso e distribuzione, ma significa accettare una perdita reale sulla fluidità vocale rispetto all'esecuzione desktop nativa. Per lo stesso motivo, nella versione attuale si è preferito disabilitare in WebGL il \emph{reply word flow} temporizzato e ricorrere a una visualizzazione testuale statica del reply: è una scelta di stabilità, non un dettaglio lasciato indietro. La resa del lipsync è ancora sensibile all'integrazione concreta tra Avaturn, \texttt{uLipSync} e condizioni WebGL. Anche con fallback e adattamenti sul mapping dei blendshape, nel prototipo non è stato possibile ottenere una sincronizzazione uniforme in tutte le combinazioni di avatar e carico. In pratica il problema WebGL non è l'assenza di voce, ma la difficoltà di sostenere in modo robusto una sincronizzazione fine tra audio, testo e movimento labiale. Portare questo livello a maturità richiederebbe una riscrittura più adatta del layer di lipsync e sincronizzazione per la pipeline WebGL.

Tra i limiti emersi c'è la tendenza del modello linguistico a colmare vuoti di contesto con contenuto plausibile quando il recupero delle fonti non porta materiale davvero utile. In \autoref{subsec:criticita-osservate} questo comportamento è descritto dal punto di vista osservativo; qui pesa come limite strutturale del bilanciamento tra fluidità conversazionale e controllo semantico. In un contesto embodied, anche una deviazione plausibile pesa più di un errore testuale isolato perché viene attribuita all'identità dell'interlocutore. Mitigare questo problema richiede controlli aggiuntivi sull'aderenza ai contenuti recuperati o politiche più conservative quando il contesto è sottosoglia, ma entrambe aumentano latenza o complessità nel perimetro attuale. Quando i PDF ingestiti contengono tabelle articolate, layout multi-colonna o scansioni a bassa qualità, l'OCR produce testo frammentato e spesso riordinato in modo non lineare. Quei chunk rumorosi finiscono comunque nello store vettoriale e possono essere recuperati con similarità alta anche se il contenuto è parziale o invertito. Il sistema continua quindi a produrre una risposta, ma il contesto che costruisce può essere poco utilizzabile per la generazione, con esiti confusi o dichiarazioni di incertezza anche su documenti presenti. Qui si vede bene il limite della scelta pragmatica adottata: l'ingestione OCR generalista rende il sistema subito utilizzabile su più formati, ma lascia entrare rumore strutturale che il retrieval ibrido da solo non riesce sempre a correggere. Per affrontare davvero questo punto serve un pre-processing \textit{layout-aware} prima dell'indicizzazione, che non è stato ancora integrato.

Anche la qualità del voice cloning presenta un compromesso. Coqui XTTS v2 produce una voce riconoscibile come sintetica e, su frasi lunghe o punteggiatura complessa, timbro e prosodia possono allontanarsi dal campione registrato. La causa principale è la combinazione tra campione breve raccolto in app e vincoli computazionali della macchina di sviluppo: estendere durata e qualità del riferimento migliorerebbe la resa, ma allungherebbe onboarding e inferenza oltre un equilibrio sostenibile sul laptop usato per i test (AMD Ryzen AI 9 HX 370 con RTX 5070 Ti Mobile, 12 GB di VRAM). Pre-attivazione, streaming e \emph{wait phrases} attenuano la percezione dell'attesa, ma non cambiano questo compromesso di fondo: la risposta può partire prima, non diventare improvvisamente più fedele al riferimento o meno costosa da generare. Per il perimetro attuale, XTTS v2 è la scelta più praticabile per voice cloning zero-shot locale, ma la somiglianza percepita va trattata come approssimazione e non come fedeltà.

Infine c'è un limite operativo che non va nascosto: la toolchain Python/CUDA resta fragile quando cambiano dipendenze o driver. Dal punto di vista infrastrutturale il sistema ha già un perimetro minimo abbastanza chiaro, perché in deploy i servizi restano sull'indirizzo locale della macchina e il traffico pubblico passa per reverse proxy e TLS; manca però ancora un vero irrobustimento applicativo, con autenticazione, autorizzazione e limitazione del traffico coerenti con un'esposizione più ampia. Resta anche un limite sulla privacy operativa. Campioni vocali, note personali, PDF, immagini descritte e profili avatar possono contenere dati sensibili o comunque identificativi; questi materiali sono trattati assumendo un ambiente controllato e un uso sperimentale. Finché autenticazione, politiche di accesso, cifratura applicativa e gestione esplicita del consenso rimangono parziali, SOULFRAME non è adatto a un impiego aperto o multiutente con dati personali non presidiati.

\section{Sviluppi futuri prioritari}
\label{sec:sviluppi-futuri}
\paragraph{Miglioramenti tecnici.}
Uno sviluppo prioritario riguarda la memoria conversazionale. Non si tratta solo di superare il trigger esplicito, ma di introdurre un modulo che estragga automaticamente le informazioni rilevanti dal dialogo (ad esempio NER\footnote{NER sta per \emph{Named Entity Recognition}: è il riconoscimento automatico di entità nominate nel testo, come persone, luoghi, date o organizzazioni. In un sistema come SOULFRAME potrebbe aiutare a capire quali elementi del dialogo valgano la pena di essere memorizzati.}, intent detection o un secondo LLM leggero in pipeline), mantenendo però regole conservative per evitare memorizzazioni spurie. Sul solo sottosistema RAG avrebbe senso affiancare a questo hardening una batteria controllata e ripetibile lato backend, con telemetria per turno, routing dell'intent più robusto e controlli più prudenti quando il contesto recuperato è debole o rumoroso. In parallelo, varrebbe la pena provare modelli linguistici più recenti e più capaci, perché una parte dei limiti osservati nella risposta finale dipende anche dalla taglia e dalla qualità del modello oggi usato. 

Un secondo fronte riguarda la resa vocale. Nel prototipo attuale il voice cloning riesce spesso a catturare bene l'andamento generale della voce, ma resta meno convincente proprio nei campioni più deboli o meno marcati, dove la somiglianza percepita cala in modo evidente. Una possibile evoluzione è sperimentare modelli TTS più pesanti o più aggiornati, anche a costo di un aumento della latenza, per capire se sia possibile migliorare fedeltà timbrica, tenuta delle voci meno ``forti'' e naturalezza complessiva del parlato.

Resta aperto anche il tema di una personalizzazione avatar più spinta a partire da immagini reali. Non avrebbe senso aggiungere complessità per principio: questa direzione varrebbe la pena solo se il salto percettivo fosse concreto e sostenibile rispetto alla qualità già garantita dal flusso attuale. Allo stesso modo, un miglioramento interessante del lato conversazionale sarebbe mostrare in interfaccia le fonti che l'avatar sta richiamando mentre parla, ad esempio un'immagine, un estratto di PDF o un documento recuperato dalla memoria. Questo renderebbe più trasparente il legame tra risposta e contenuti usati dal backend, e aiuterebbe anche a capire meglio quando la risposta finale sta realmente seguendo il contesto disponibile.

Anche un'estensione verso mobile/Android ha senso, soprattutto perché l'interfaccia è già impostata con logica touch-first in vari passaggi, mentre l'inferenza rimane legata a un backend dedicato. Un'architettura ibrida (client mobile + modelli lato cloud) potrebbe abbassare la barriera d'accesso senza comprimere eccessivamente la qualità. Prima di allargare troppo il perimetro, serve però consolidare la base tecnica: telemetria per turno, logging strutturato, readiness-check più rigorosi e hardening operativo minimo del deploy. Sono interventi meno ``visibili'', ma necessari per trasformare le osservazioni qualitative in evidenze misurabili. Serve capire meglio se il limite stia nel retrieval, nella selezione del contesto o nella risposta finale.

\paragraph{Estensione della valutazione.}
Nel perimetro di questa tesi la valutazione rimane centrata su osservazioni tecniche e qualitative del prototipo in esercizio. In prospettiva, il passo successivo più utile è introdurre protocolli ripetibili, telemetria per turno e un confronto più controllato tra ambienti, così da trasformare le osservazioni esplorative in evidenze più robuste.

\enlargethispage{2\baselineskip}
\label{sec:finali}
SOULFRAME mostra che una conversazione vocale embodied in WebGL diventa praticabile solo quando architettura, implementazione e operatività vengono tenute insieme fin dall'inizio. Il contributo principale del lavoro sta nell'aver trasformato componenti eterogenei in un prototipo unico, funzionante e abbastanza trasparente da poter essere corretto, misurato e sviluppato con criterio.
